% !TeX root = ../../tesis.tex

% ============================================================
%  §5.6.1 — Contexto del instrumento para el diagnóstico
% ============================================================

Los cinco indicadores derivados de las Fases~1 a~3 ---$C$, $C_i$, $DI$, $T$ y
$B(v)$--- caracterizan la red desde dimensiones que no comparten unidad de
medida. La \textbf{Clasificación Garibelt}, definida formalmente en la
\autoref{subsec:clasificacion-garibelt}, es el instrumento que reúne esas cinco
medidas en un tablero de diagnóstico topológico sin reducirlas a un escalar
único: cada dimensión conserva su rango, su método de normalización y su
significado estructural propio (\autoref{tab:clasificacion-garibelt-dims}).

El procedimiento opera en dos pasos. Primero, cada valor bruto se normaliza a
$[0,\,1]$ con el método que corresponde a sus propiedades estadísticas:
referencia externa para $T$, límite teórico para $DI$, percentil empírico para
$C_i$ e índice de Gini invertido para $B(v)$~(
\autoref{tab:clasificacion-garibelt-dims}). Segundo, el valor normalizado se
clasifica en una de las cuatro bandas de la \textbf{Escala Garibelt}: Crítico
$[0{,}00,\,0{,}25)$, Débil $[0{,}25,\,0{,}50)$, Aceptable $[0{,}50,\,0{,}75)$ e
Idóneo $[0{,}75,\,1{,}00]$. El resultado de aplicar este procedimiento a los
valores calculados en las Fases~1 a~3 es el \textbf{Espectro MJG}: el perfil de
cinco bandas que expone, sin suavizarlas, las fortalezas y las vulnerabilidades
estructurales de la red de la \zmvm{}.
